\documentclass[11pt,a4paper]{article}

\usepackage[utf8]{inputenc}
\usepackage[T1]{fontenc}
\usepackage[spanish,es-tabla,es-nodecimaldot]{babel}
\usepackage{lmodern}
\usepackage[margin=2.5cm]{geometry}
\usepackage{graphicx}
\usepackage{xcolor}
\usepackage{booktabs}
\usepackage{tabularx}
\usepackage{longtable}
\usepackage{array}
\usepackage{enumitem}
\usepackage{amssymb}
\usepackage{float}
\usepackage{multicol}
\usepackage{pdfpages}
\usepackage{tikz}
\usetikzlibrary{positioning,arrows.meta,shapes.geometric}
\usepackage{fancyhdr}
\usepackage[font=small,labelfont=bf]{caption}
\usepackage{url}
\usepackage[hidelinks]{hyperref}
\hypersetup{
  pdftitle={TA-PFC-E1 - UteqFit-AI: visión, alcance, partes interesadas, contexto y plan de elicitación},
  pdfauthor={Aldas Zambrano A., Briones Solorzano W., Lavayen Bone C., Monserrate Ochoa E.},
  pdfsubject={Ingeniería de Requerimientos (ISR-401) - UTEQ},
  pdfkeywords={ingeniería de requisitos, elicitación, partes interesadas, UteqFit-AI}
}

% Encabezado y pie de página
\pagestyle{fancy}
\fancyhf{}
\fancyhead[L]{\small ISR-401 \textperiodcentered{} TA-PFC-E1}
\fancyhead[R]{\small UteqFit-AI}
\fancyfoot[C]{\small \thepage}
\renewcommand{\headrulewidth}{0.4pt}
\setlength{\headheight}{14pt}
\fancypagestyle{plain}{\fancyhf{}\fancyfoot[C]{\small \thepage}\renewcommand{\headrulewidth}{0pt}}

% Evita líneas que se salen del margen
\setlength{\emergencystretch}{2em}

% Rutas: compila igual desde PFC/E1 (local) o desde la raíz del proyecto (Overleaf)
\IfFileExists{PFC/E1/main.tex}{\def\rutaE{PFC/E1/}}{\def\rutaE{}}
\makeatletter
\def\input@path{{\rutaE}}
\makeatother
\graphicspath{{\rutaE}{\rutaE figuras/}}

\newcolumntype{L}{>{\raggedright\arraybackslash}X}

% Marca de contenido por completar. Antes de la entrega final no debe quedar ninguna:
%   grep -rn "pendiente" PFC/E1
\newcommand{\pendiente}[1]{\textcolor{red}{[PENDIENTE: #1]}}

% Incluye un archivo si ya existe; si no, deja un aviso y el documento compila igual
\newcommand{\incluir}[1]{\IfFileExists{\rutaE #1.tex}{\input{#1}}{\pendiente{falta subir el archivo \texttt{\detokenize{#1}.tex}}}}

\newcommand{\sistema}{UteqFit-AI}
\newcommand{\organizacion}{Gimnasio Universitario de la Unidad de Bienestar Universitario (UBU) de la UTEQ}
\newcommand{\repourl}{https://github.com/AlanAldasZam/UteqFit-AI}

\begin{document}

% ---------------- Portada ----------------
\begin{titlepage}
  \centering
  {\large\bfseries Universidad Técnica Estatal de Quevedo\par}
  {Facultad de Ciencias de la Computación\par}
  {Carrera de Software\par}
  \vspace{2.5cm}
  {Ingeniería de Requerimientos (ISR-401)\par}
  {4.\textsuperscript{o} nivel \textperiodcentered{} Paralelo A \textperiodcentered{} Período SPA 2026--2027\par}
  \vspace{2.5cm}
  {\Large\bfseries Proyecto Fin de Curso -- Entrega 1\par}
  \vspace{0.3cm}
  {\large Visión, alcance, partes interesadas, contexto y plan de elicitación\par}
  \vspace{0.3cm}
  {\large Proyecto: \sistema\par}
  \vspace{2cm}
  {\bfseries Integrantes\par}
  Aldas Zambrano Alan Andrés\par
  Briones Solorzano William Adrian\par
  Lavayen Bone Cristopher Javier\par
  Monserrate Ochoa Erick Ariel\par
  \vspace{1.5cm}
  {\bfseries Docente\par}
  Gleiston Cicerón Guerrero Ulloa, PhD\par
  \vfill
  Repositorio: \url{\repourl}\par
  \vspace{0.3cm}
  Código de la tarea: TA-PFC-E1\par
  Quevedo, noviembre de 2026\par
\end{titlepage}

\tableofcontents
\listoftables
\listoffigures
\clearpage

% ---------------- Cuerpo (7 secciones de la guía) ----------------
\section{Organización real}\label{sec:organizacion}
\incluir{secciones/01_organizacion}

\section{Visión y alcance}\label{sec:vision}
\incluir{secciones/02_vision_alcance}

\section{Partes interesadas}\label{sec:partes}
\incluir{secciones/03_partes_interesadas}

\section{Contexto}\label{sec:contexto}
\incluir{secciones/04_contexto}

\section{Plan de elicitación}\label{sec:elicitacion}
\incluir{secciones/05_plan_elicitacion}

\section{Ética y datos}\label{sec:etica}
\incluir{secciones/06_etica_datos}

\section{Gestión del proyecto}\label{sec:gestion}
\incluir{secciones/07_gestion_proyecto}
\incluir{secciones/08_riesgos}

\section*{Declaración de uso de inteligencia artificial}\label{sec:ia}
\addcontentsline{toc}{section}{Declaración de uso de inteligencia artificial}
\incluir{secciones/09_declaracion_ia}

\bibliographystyle{IEEEtran}
\bibliography{\rutaE referencias}

% ---------------- Anexos ----------------
\clearpage
\appendix
\section{Guion de entrevista}\label{anx:entrevista}
\incluir{instrumentos/guion_entrevista}

\section{Cuestionario}\label{anx:cuestionario}
\incluir{instrumentos/cuestionario}

\section{Guía de observación}\label{anx:observacion}
\incluir{instrumentos/guia_observacion}

% El consentimiento se inserta como PDF cuando ya fue compilado en instrumentos/
\IfFileExists{\rutaE instrumentos/Consentimiento_Informado_PFC_ISR401.pdf}{%
  \includepdf[pages=-, pagecommand={\thispagestyle{empty}}, addtotoc={1,section,1,Modelo de consentimiento informado,anx:consentimiento}]{\rutaE instrumentos/Consentimiento_Informado_PFC_ISR401.pdf}%
}{%
  \section{Modelo de consentimiento informado}\label{anx:consentimiento}
  \pendiente{compilar instrumentos/Consentimiento\_Informado\_PFC\_ISR401.tex}%
}

\end{document}